\subsection{Robust compilation from hidden targets}\label{ian:robust-compilation}

Status: proved: distributional quantum readout compiles to a globally defined acceptance function without additional queries or loss of sector weight; the variance and covariance perturbation constants are sharp as universal linear constants. Uses Definition~\ref{ian:hard-sector-definition}, Proposition~\ref{ian:hard-sector-forward} and Lemma~\ref{ian:hard-sector-aa-exclusion} of \path{ian/13_hard_sector_definitions_forward.tex}, Remark~\ref{ian:hard-sector-polynomial-quantifiers} of \path{ian/15_hard_sector_sets_asymptotics.tex}, and Lemma~\ref{ian:hard-sector-orbit-covariance} of \path{ian/17_hard_sector_instances.tex}; the Bernstein--Vazirani readout in Proposition~\ref{ian:robust-compilation-character-check} is cited from the literature and its action is verified in place.

\begin{definition}[Joint law, oracle, and decoder error]
\label{ian:robust-compilation-conventions}
Let \(\Omega_n=\{-1,1\}^n\), let \(\mu\) be an ambient probability law, and let \(\pi\) be a specified joint law of
\[
 (Z,X)\in\mathcal L\times\Omega_n,
 \qquad X\sim\nu,
\]
where \(\mathcal L\) is finite. Fix \(g:\mathcal L\to[0,1]\) and put \(Y=g(Z)\).
A deterministic label means \(Z=\lambda(X)\) almost surely; a latent label need not be a function of \(X\).
Expectations involving \(Z\) or \(Y\) are under \(\pi\); expectations of functions of \(X\) alone are under \(\nu\).
Thus notation such as \(\E_\nu|Q-Y|\), when used for a latent target, means \(\E_\pi|Q(X)-Y|\).

The sector law has weight \(w\in(0,1]\) if either
\[
 w\nu(x)\le\mu(x)\quad\text{for every }x
\]
in the smooth interpretation, or
\[
 \nu=\mu(\cdot\mid H),\qquad \mu(H)\ge w
\]
for a specified set \(H\) in the set interpretation.
No hardness assertion is included in these preliminary conditions.

Coordinate access means the standard oracle
\[
 O_x|i,b\rangle
 =|i,b\oplus(1-x_i)/2\rangle,
 \qquad i\in\{1,\ldots,n\},\quad b\in\{0,1\}.
\]
Input-independent operations and postprocessing are free. Phase access follows by applying this oracle with the answer register in
\(\lvert-\rangle=(\lvert0\rangle-\lvert1\rangle)/\sqrt2\).
There is no additional access to \(Z\), sector membership, or a sampler.

A \(q\)-query decoder is one fixed circuit, defined for every \(x\in\Omega_n\), using at most \(q\) oracle calls.
Write \(K_x\) for its output law on input \(x\).
Its internal randomness is independent of \(Z\) conditional on \(X\), so the output's conditional law given \((X,Z)\) is \(K_X\).
For a label output \(\widehat Z\), distributional label error at most \(\delta\) means
\[
 \Pr_{\pi,\mathrm{decoder}}[\widehat Z\ne Z]\le\delta.
\]
More generally, a bounded target output \(\widehat Y\in[0,1]\) has distributional target error at most \(\delta\) if
\[
 \E_{\pi,\mathrm{decoder}}|\widehat Y-g(Z)|\le\delta.
\]
Label error implies target error for \(\widehat Y=g(\widehat Z)\).
A failure symbol is assigned a fixed value in \([0,1]\) and counts as a label error.
\end{definition}

Throughout, \(\mathcal D_t\) and \(\mathcal F_t\) are deterministic coordinate trees of worst-case depth at most \(t\), with leaf outputs in \([0,1]\) and \([-1,1]\), respectively, as in
\path{ian/13_hard_sector_definitions_forward.tex}.
In particular,
\[
 \mathcal F_t=\{2D-1:D\in\mathcal D_t\}.
\]

\begin{lemma}[Bounded-variable perturbation]
\label{ian:robust-compilation-perturbation}
Let \(U,V\in[0,1]\) be random variables on the same probability space, with
\(\E|U-V|\le\delta\). Then
\[
 |\Var U-\Var V|\le2\delta.
\]
For every random variable \(D\in[0,1]\), and every \(F\in[-1,1]\),
\[
 |\Cov(U,D)-\Cov(V,D)|\le\delta,
 \qquad
 |\Cov(U,F)-\Cov(V,F)|\le2\delta.
\]
\end{lemma}

\begin{proof}
For \(u,v,c\in[0,1]\),
\[
 |(u-c)^2-(v-c)^2|
 =|u-v|\,|u+v-2c|
 \le2|u-v|.
\]
Using \(c=\E U\) gives
\[
 \Var V
 \le\E(V-\E U)^2
 \le\Var U+2\delta.
\]
Interchanging \(U,V\) proves the variance assertion.

Centering only the test gives
\[
 \Cov(U,D)-\Cov(V,D)
 =\E[(U-V)(D-\E D)].
\]
Since \(|D-\E D|\le1\), this has absolute value at most \(\delta\).
The same argument uses \(|F-\E F|\le2\) for the signed test.
\end{proof}

\begin{theorem}[Robust compilation]
\label{ian:robust-compilation-theorem}
Use Definition~\ref{ian:robust-compilation-conventions}, with integers \(t,q\ge0\).
Suppose \(0<v\le1/4\), \(\eta\ge0\), \(0\le\delta\le1\), and
\[
 \Var_\pi Y\ge v,
 \qquad
 \sup_{F\in\mathcal F_t}
 |\Cov_\pi(Y,F(X))|\le\eta.
 \tag{28.1}\label{ian:robust-compilation-LH}
\]
Suppose a \(q\)-query decoder produces a bounded target output \(\widehat Y\) with distributional target error at most \(\delta\).
This includes a label decoder with distributional label error at most \(\delta\), followed by \(g\).

Running the decoder and accepting with probability \(\widehat Y\) produces a globally defined acceptance function \(Q:\Omega_n\to[0,1]\), using
\[
 T=q
\]
queries, such that
\[
 \E_\pi|Q(X)-Y|\le\delta,
 \qquad
 \Var_\nu Q\ge v-2\delta,
\]
and, for every \(D\in\mathcal D_t\) and \(F\in\mathcal F_t\),
\[
 |\Cov_\nu(Q,D)|\le\frac{\eta}{2}+\delta,
 \qquad
 |\Cov_\nu(Q,F)|\le\eta+2\delta.
 \tag{28.2}\label{ian:robust-compilation-LR}
\]
For \(q\ge1\), \(Q\in\mathcal Q_{n,q}\).

Put
\[
 a=v-2\delta,\qquad \gamma_0=\eta+2\delta.
\]
For every positive \(\gamma\) satisfying \(\gamma_0\le\gamma<a\),
\[
 \operatorname{Hard}(Q,\nu;t,a,\gamma).
\]
The original set or smooth sector retains weight \(w\).

More generally, the same variance, covariance, hardness, and weight conclusions hold for every supplied globally defined \(Q:\Omega_n\to[0,1]\) satisfying only
\[
 \E_\pi|Q(X)-Y|\le\delta.
\]
In this general version, if \(Q\) is the acceptance function of a quantum algorithm using at most \(T\ge1\) coordinate queries, then \(Q\in\mathcal Q_{n,T}\) and \(\operatorname{Hard}(Q,\nu;t,a,\gamma)\) for every such \(\gamma\).
\end{theorem}

\begin{proof}
Let \(K_x\) be the decoder's output law on input \(x\), and define
\[
 Q(x)=\E_{K_x}\widehat Y.
\]
For label output this is
\[
 Q(x)=\sum_{\ell}K_x(\ell)g(\ell),
\]
including the assigned value of any failure symbol.
Measuring the decoder output and tossing the indicated acceptance coin uses no further oracle calls.
The fixed circuit and postprocessing define \(Q\) on the entire cube, including inputs outside the support of \(\nu\).

Since the output's conditional law given \((X,Z)\) is \(K_X\),
\[
 \E[\widehat Y\mid X,Z]=Q(X).
\]
Consequently, conditional Jensen yields
\[
 \E_\pi|Q(X)-Y|
 =\E\bigl|\E[\widehat Y-Y\mid X,Z]\bigr|
 \le\E|\widehat Y-Y|
 \le\delta.
\]
For a label decoder, the last inequality also follows from
\[
 |g(\widehat Z)-g(Z)|
 \le\mathbf1_{\{\widehat Z\ne Z\}}.
\]
This argument applies to latent labels without identifying \(Y\) with a function of \(X\).

Apply Lemma~\ref{ian:robust-compilation-perturbation} on the joint probability space to \(U=Q(X)\), \(V=Y\).
The hypothesis for signed trees implies
\[
 |\Cov_\pi(Y,D(X))|\le\eta/2,
\]
because \(2D-1\in\mathcal F_t\).
The claimed variance and covariance estimates follow.

If \(0<\gamma<a\) and \(\gamma\ge\gamma_0\), then \(0<a\le v\le1/4\), and these estimates are precisely
Definition~\ref{ian:hard-sector-definition} in
\path{ian/13_hard_sector_definitions_forward.tex}.

The law \(\nu\) and its domination inequality have not changed.
For a set sector, the same \(H\) still has \(\mu(H)\ge w\) and conditional law \(\nu\); also
\[
 w\nu(x)
 =\frac{w}{\mu(H)}\mu(x)\mathbf1_H(x)
 \le\mu(x).
\]
For a smooth sector, the original inequality \(w\nu\le\mu\) already supplies the required weight.
The proof after the construction of \(Q\) uses only the displayed \(L^1\) approximation, establishing the general version.
\end{proof}

\begin{corollary}[Exact certificate regimes and forward loss]
\label{ian:robust-compilation-regimes}
Under Theorem~\ref{ian:robust-compilation-theorem}, retain
\[
 a=v-2\delta,\qquad \gamma_0=\eta+2\delta.
\]
The canonical parameters satisfy \(0<\gamma_0<a\le1/4\) exactly when
\[
 \eta+2\delta>0,\qquad \eta+4\delta<v.
\]
A tolerance \(\gamma\) with \(\gamma>0\) and \(\gamma_0\le\gamma<a\) exists exactly when
\[
 \eta+4\delta<v.
\]
A tolerance \(\gamma\) with \(\gamma>0\) and \(\gamma_0\le\gamma\le a/4\) exists exactly when
\[
 4\eta+10\delta\le v.
 \tag{28.3}\label{ian:robust-compilation-quarter-condition}
\]
If \(\gamma_0>0\), this last inequality is exactly \(\gamma_0\le a/4\).
If \(\eta=\delta=0\), take \(a=v\) and any \(0<\gamma\le v/4\), for example \(\gamma=v/4\).

Whenever \(\eta+4\delta<v\), the loss satisfies
\[
 \ell_\mu(t,Q)\ge w(v-\eta-4\delta).
\]
In the quarter-tolerance regime,
\[
 \operatorname{Hard}(Q,\nu;t,a,a/4),
 \qquad
 \ell_\mu(t,Q)\ge\frac34w(v-2\delta).
\]
If full AA holds with constants \(A_0,B_0\), then in this regime
\[
 t<
 \left\lceil
 A_0\left(T_*+\frac{8}{3w(v-2\delta)}\right)^{B_0}
 \right\rceil,
 \qquad T_*=\max\{1,q\}.
\]
For a supplied \(T\)-query acceptance function, replace \(T_*\) by \(\max\{1,T\}\).
\end{corollary}

\begin{proof}
The strict inequality \(\gamma_0<a\) is
\(\eta+4\delta<v\).
The positivity requirement excludes only \(\eta=\delta=0\) for the canonical tolerance.
Moreover,
\[
 \gamma_0\le a/4
 \quad\Longleftrightarrow\quad
 4\eta+10\delta\le v.
\]
Since \(v>0\), this last inequality ensures \(a>0\).
When \(\gamma_0=0\), choosing \(a/4\) supplies a positive tolerance.

For \(\gamma_0>0\), Proposition~\ref{ian:hard-sector-forward} in
\path{ian/13_hard_sector_definitions_forward.tex} gives
\[
 \ell_\mu(t,Q)
 \ge w(a-\gamma_0)
 =w(v-\eta-4\delta).
\]
If \(\gamma_0=0\), apply that proposition with every \(\gamma\in(0,a)\) and let \(\gamma\) decrease to zero.
In the quarter regime, increase the tolerance to \(a/4\) and apply the same proposition.

The AA depth bound is Lemma~\ref{ian:hard-sector-aa-exclusion} with this value of \(a\).
The maximum with \(1\) only accommodates the convention that
\(\mathcal Q_{n,T}\) in \eqref{ian:hard-sector-aa-definition} is indexed by \(T\ge1\).
\end{proof}

\begin{remark}[Covariance normalization and asymptotic quantifiers]
\label{ian:robust-compilation-normalization}
The parameter \(\eta\) above uses the signed-tree normalization of \eqref{ian:robust-compilation-LH}.
If the hiding hypothesis is instead
\[
 \sup_{D\in\mathcal D_t}
 |\Cov_\pi(Y,D(X))|\le\eta_D,
\]
substitute \(\eta=2\eta_D\). In particular,
\[
 |\Cov_\nu(Q,D)|\le\eta_D+\delta,
 \qquad
 a=v-2\delta,\qquad
 \gamma_0=2\eta_D+2\delta.
\]
The corresponding strict and quarter-tolerance criteria are
\[
 2\eta_D+4\delta<v,
 \qquad
 8\eta_D+10\delta\le v,
\]
with the positive-tolerance convention of
Corollary~\ref{ian:robust-compilation-regimes}.

For an asymptotic full-AA failure conclusion, compiled witnesses must satisfy the quarter-tolerance condition \eqref{ian:robust-compilation-quarter-condition}, or its \(\eta_D\) form above, so that \(\operatorname{Hard}(Q,\nu;t,a,a/4)\) holds with \(a=v-2\delta\).
They must also satisfy the all-budget requirements of \eqref{ian:hard-sector-polynomial-property}: there are \(K,p,p'\ge1\) such that for every \(A,B\ge1\) some compiled witness has
\[
 r>8,
 \qquad
 T=\max\{1,q\},
 \qquad
 t>\left\lceil A(T+r)^B\right\rceil,
 \qquad
 w^{-1}\le Kr^p,
 \qquad
 (v-2\delta)^{-1}\le Kr^{p'},
\]
that is, with \(w^{-1}\) and \((v-2\delta)^{-1}\) bounded by fixed polynomials in \(r\).
The polynomial bounds are chosen before the budget quantifiers, as specified in
Remark~\ref{ian:hard-sector-polynomial-quantifiers} of
\path{ian/15_hard_sector_sets_asymptotics.tex} and
Theorem~\ref{ian:hard-sector-aa-equivalence} of
\path{ian/13_hard_sector_definitions_forward.tex}.
\end{remark}

\begin{proposition}[Sharp perturbation constants]
\label{ian:robust-compilation-sharpness}
The coefficients \(2\) in the variance bound, \(1\) in the \(\mathcal D_t\) covariance bound, and \(2\) in the \(\mathcal F_t\) covariance bound cannot be decreased as universal linear coefficients.
These assertions hold even for deterministic labels, exact hiding at depth one, and decoders of bounded query cost.
\end{proposition}

\begin{proof}
For variance, let \(B\) be Bernoulli with parameter \(p\in(0,1)\), let \(R\) be an independent uniform sign, and set
\[
 X=(R,R(2B-1)),\qquad Y=B,\qquad
 Q(x)=(1-\epsilon)\frac{1+x_1x_2}{2},
 \qquad 0<\epsilon<1.
\]
Each coordinate is independent of \(B\), so \(Y\) has zero covariance with every depth-one tree.
A decoder reads both coordinates, computes \(B\), and changes a recovered \(1\) to \(0\) with probability \(\epsilon\).
With \(g(b)=b\), its acceptance function is \(Q\), and
\[
 \delta=\E|Q(X)-Y|=p\epsilon.
\]
Its variance loss is
\[
 \Var Y-\Var Q
 =(2\epsilon-\epsilon^2)p(1-p)
 =\delta(2-\epsilon)(1-p).
\]
Taking \(p=\epsilon\) and letting \(p\) decrease to zero makes the ratio of variance loss to \(\delta\) tend to \(2\).
In this family, the coefficient \(2\) is attained only as \(p\to0\): at fixed \(p,\epsilon\in(0,1)\) the ratio \((2-\epsilon)(1-p)\) is strictly below \(2\).

For covariance, let \(X_1,X_2\) be independent uniform signs, and let \(X_3\) be independent with \(\Pr[X_3=1]=p\).
Put
\[
 Y=\frac{1+X_1X_2}{2},\qquad
 D=\frac{1+X_3}{2},\qquad
 Q=Y+(1-Y)D.
\]
Again \(Y\) has zero covariance with every depth-one tree, and \(\Var Y=1/4\).
A three-query decoder computes these bits and outputs \(Y\lor D\); its label error against \(Y\) is
\[
 \delta=\Pr[Y=0,D=1]=p/2.
\]
Here
\[
 \Cov(Q,D)=\frac{p(1-p)}2=\delta(1-p),
 \qquad
 \Cov(Q,2D-1)=2\delta(1-p).
\]
Letting \(p\) decrease to zero proves both covariance assertions.
These examples establish optimal universal linear coefficients, without asserting an optimal nonlinear bound at each fixed positive \(\delta\).
\end{proof}

\begin{remark}[The composed forward-loss error is not sharp]
\label{ian:robust-compilation-direct-loss}
Sharpness of the separate perturbation constants does not make the composed loss \(w(v-\eta-4\delta)\) sharp.
Under just the target hiding and \(L^1\) approximation hypotheses, a direct argument gives
\[
 \ell_\mu(t,Q)\ge w\max\{0,v-\eta-\delta\}.
\]
Indeed, for \(D\in\mathcal D_t\), with \(m=\E_\pi Y\),
\[
 \E_\pi(Y-D(X))^2
 =\Var_\pi Y+\E_\nu(D-m)^2
   -2\Cov_\pi(Y,D(X))
 \ge v-\eta.
\]
Since \(Y,D(X)\in[0,1]\), one has \(|Y-D(X)|\ge(Y-D(X))^2\), so \(\E_\pi|Y-D(X)|\ge v-\eta\); the triangle inequality then gives
\[
 \E_\nu|Q-D|
 \ge\E_\pi|Y-D(X)|-\delta
 \ge v-\eta-\delta.
\]
Multiply by \(w\) using \(\mu\ge w\nu\), minimize over \(D\), and use nonnegativity.
The triangle-inequality step is the comparison \eqref{ian:phase-label-pure-target-comparison} of
Proposition~\ref{ian:phase-label-flat-components} in
\path{ian/20_phase_label_theorem.tex}, with the target \(Y\) in the role of \(g(Z)\) and \(e=\delta\), applied with \(\nu\) (joint law \(\pi\)) in place of \(\mu\), before minimizing over \(D\).
The hard-sector certificate remains useful because it supplies the variance--covariance hypotheses required by the AA exclusion and equivalence results.
\end{remark}

\begin{proposition}[Character-code sanity check]
\label{ian:robust-compilation-character-check}
Let \(m\ge2\) be even, \(n=2^m\), and identify coordinates with \(z\in\mathbb F_2^m\).
Let
\[
 c_a(z)=(-1)^{a\cdot z},\qquad
 C_n=\{c_a:a\in\mathbb F_2^m\},
 \qquad A\sim U_{\mathbb F_2^m},\qquad X=c_A.
\]
Take the deterministic label \(Z=A\) and the quadratic target
\[
 Y=h(A),\qquad
 h(a)=\sum_{j=1}^{m/2}a_{2j-1}a_{2j}\pmod2.
\]
There is a one-query decoder with \(\delta=0\).
Compilation gives the globally defined one-query acceptance function
\[
 Q(x)=P(x)
 =\sum_{a:h(a)=1}\widehat x(a)^2,
 \qquad
 \widehat x(a)=\frac1n\sum_zx_z(-1)^{a\cdot z}.
\]
The target parameters are
\[
 v=\frac{1-1/n}{4},
 \qquad
 \eta=\frac{2^t-1}{2\sqrt n}.
\]
Consequently, for \(n\ge16\) and \(t\le m/2-2\),
\[
 \operatorname{Hard}\left(
 P,U_{C_n};t,\frac3{16},\frac18\right),
\]
recovering Proposition~\ref{ian:hard-sector-character-instance} in
\path{ian/17_hard_sector_instances.tex}.
Every ambient set or smooth realization of \(U_{C_n}\) retains its original weight.
\end{proposition}

\begin{proof}
The decoder is Bernstein--Vazirani readout (from literature: Ethan Bernstein and Umesh Vazirani, \emph{Quantum Complexity Theory}, SIAM Journal on Computing \textbf{26}(5), 1411--1473 (1997), DOI \texttt{10.1137/S0097539796300921}).
Only exact character-oracle recovery is used; its action is verified here.

Prepare the uniform address state, apply one coordinate query by phase kickback, and apply the Walsh transform.
For every cube input this gives
\[
 \frac1{\sqrt n}\sum_z|z\rangle
 \ \longmapsto\
 \frac1{\sqrt n}\sum_zx_z|z\rangle
 \ \longmapsto\
 \sum_a\widehat x(a)|a\rangle.
\]
For \(x=c_A\), character orthogonality makes the final state exactly \(|A\rangle\).
Thus \(q=T=1\) and \(\delta=0\).
Measuring \(a\) and accepting when \(h(a)=1\) gives \(P\).
Parseval gives
\[
 \sum_a\widehat x(a)^2=\frac1n\sum_zx_z^2=1,
\]
so \(P\in[0,1]\) globally.

For a uniform pair of bits, the mean of \((-1)^{a_1a_2}\) is \(1/2\).
Independence across the \(m/2\) pairs therefore gives
\[
 \E(-1)^{h(A)}=n^{-1/2},\qquad
 \Pr[h(A)=1]=\frac{1-n^{-1/2}}2,
 \qquad
 \Var h(A)=\frac{1-1/n}{4}.
\]
Since the final state on \(c_A\) is exactly \(|A\rangle\), the compiled function satisfies \(P(c_A)=h(A)\); thus \(P(X)=h(A)\) under \(U_{C_n}\).
Lemma~\ref{ian:hard-sector-orbit-covariance} in
\path{ian/17_hard_sector_instances.tex}, with no noise and \(k=1\), supplies
\[
 \sup_{F\in\mathcal F_t}
 |\Cov_{U_{C_n}}(h(A),F(X))|
 \le\frac{2^t-1}{2\sqrt n}.
\]
For the stated range,
\[
 \Var_{U_{C_n}}P\ge\frac3{16},
 \qquad
 \sup_{F\in\mathcal F_t}|\Cov_{U_{C_n}}(P,F)|
 \le\frac18.
\]
These are the claimed parameters, with no change of law or weight.
\end{proof}

\begin{remark}[Character weights and quarter tolerance]
\label{ian:robust-compilation-character-quarter}
For example, if
\[
 \mu=(1-\theta)U_n+\theta U_{C_n},\qquad 0\le\theta\le1,
\]
then the same set \(H=C_n\) has conditional law \(U_{C_n}\) and actual weight
\[
 \mu(C_n)=\theta+(1-\theta)n2^{-n}.
\]
The recovered coarse parameters yield
\[
 \ell_\mu(t,P)\ge\mu(C_n)/16.
\]

The coarse tolerance \(1/8\) exceeds \((3/16)/4\).
Keeping the sharper compiled values
\[
 a=\frac{1-1/n}{4},\qquad
 \gamma_0=\frac{2^t-1}{2\sqrt n},
\]
gives the quarter-tolerance certificate exactly when
\[
 2^t-1\le\frac{\sqrt n-n^{-1/2}}8.
\]
For example, \(t\le m/2-3\) suffices whenever this range is nonempty.
At \(t=0\), replace the zero tolerance by \(a/4\).
Thus the sharper certificate supplies the regime of
Lemma~\ref{ian:hard-sector-aa-exclusion} and of the hard sectors in
Remark~\ref{ian:hard-sector-polynomial-quantifiers}.
\end{remark}
